\section{行列式线丛}
\label{backdet}

我们现在将发展一些数学方法，用来帮助计算 ${\bf F}$ 型膜的赫克“本征值”.


与黎曼曲面 $D$ 上的线丛\index{line-bundle@线丛} ${\cal L}$ 相联系的是上同调群 $H^0(D,{\cal L})$ 和 $H^1(D,{\cal
L})$. ${\cal L}$ 的行列式线定义为\footnote{在代数几何中，最常采用相反的约定 $\det\,H^*({\cal L})=\left(H^0(D,{\cal
L})\right)\otimes \det\, H^1(D,{\cal L})^{-1}$，例如见 \cite{deligne}. 这里我们采用关于 $\bar\partial$ 算符行列式的文献中最常见的约定，例如见 \cite{quillen}，这也是行列式线丛\index{determinant-line-bundle@行列式线丛}进入物理学时最为人熟悉的方式. 当然，如果反转这一约定，我们的若干公式也必须相应反转.}
\begin{equation}\label{detline}\det\,H^*({\cal
L})=\det\,H^0(D,{\cal L})^{-1}\otimes \left(\det\,H^1(D,{\cal
L})\right).\end{equation} 随着 ${\cal L}$ 变化，它的行列式线也随之变化，构成 ${\rm
Pic}(D)$ 上一个称为行列式线丛的线丛 ${\bf Det}$ 的纤维. 物理学家最熟悉行列式线丛，是因为它在二维手征费米子的研究中所起的作用\transnote{二维手征费米子在背景场中的路径积分在固定背景下是一个数，但它随背景（联络或复结构）变化时并不给出函数，而给出行列式线丛的一个截面；截面上的自然度量与自然联络就是 Quillen 度量与 Quillen 联络，其曲率即费米子行列式的反常。亏格 $g$ 的黎曼曲面上，狄拉克算符族的行列式（相差一个同构）正是 theta 函数，其零除子给出 theta 除子；时空维数为 $8k+2$ 时行列式线丛还有典范平方根，即 Pfaffian 线丛（黎曼曲面正属 $k=0$ 情形）。[来源：Quillen 1985（Funct. Anal. Appl. 19，31–34，本文文献 [41]）；原文 §6；nLab "determinant line bundle"，\url{https://ncatlab.org/nlab/show/determinant+line+bundle}（Last revised 2026-01-30；2026-11-19 访问）之 "Quillen's determinant line bundle"、"Pfaffian line bundle"、"Relation to theta function" 三节。]}.

\subsection{线丛的配对}\label{linepairing}

对于我们的目的，行列式线丛的用处在于，它使我们能够定义线丛之间的某种双线性配对. 如果 ${\cal L}$ 和 ${\cal M}$ 是 $D$ 上的两个线丛，则定义如下复线 \cite{deligne}：
\begin{equation}\label{bilpa}
\langle{\cal L},{\cal M}\rangle={\det\,H^*(D,{\cal L})\otimes
\det\,H^*(D,{\cal M})\over \det\,H^*(D,{\cal L}\otimes {\cal
M})\otimes \det\,H^*(D,{\cal O})}.\end{equation} （如果 $V$ 和 $W$ 是两个一维向量空间，我们把 ${V}\otimes
{W}^{-1}$ 写成分式 $V/W$\transnote{注意式（\ref{bilpa}）右端是一维复向量空间（复线），不是复数：取行列式线即把上同调群取最高外幂，分式记号 $V/W$ 表示 $V\otimes W^{-1}$。因此“双线性”“对称”都只能理解为典范同构意义下的，这也解释了下面为何要强调 $\cong$ 可在族中同时给出。这一配对即 Deligne 的行列式—上同调配对（determinant of cohomology 配对）。[来源：Deligne 1987（Contemp. Math. 67，93–177，本文文献 [40]）；原文 §6.1；nLab "determinant line bundle"（2026-11-19 访问）。]}.）

显然存在自然同构
\begin{align}
\langle {\cal L},{\cal M}\rangle & \cong \langle {\cal M},{\cal
L}\rangle \\ \nonumber \langle {\cal L},{\cal O}\rangle \cong
\langle{\cal O},{\cal L}\rangle& \cong \Bbb{C}.\end{align} 在这里以及以下的分析中，符号 $\cong$ 指可以用一种普适方式定义、因而也能在族中给出的同构.

符号 $\langle~,~\rangle$ 最重要的性质是它的双线性，其含义是：对于任意三个线丛 ${\cal
L},\,{\cal M},$ 和 ${\cal N}$，我们都有典范同构
\begin{equation}\label{duglog} \langle{\cal L},{\cal M}\otimes {\cal 
N}\rangle\cong\langle{\cal L},{\cal M}\rangle \otimes \langle
{\cal L},{\cal N}\rangle.\end{equation} 任意线丛 ${\cal N}$ 对某个除子\index{divisor@除子} ${\DD}=\sum_in_i\qq_i$ 都具有 $\otimes_i{\cal O}(\qq_i)^{n_i}$ 的形式；因此，如果式（\ref{duglog}）对于特殊情形 ${\cal N}={\cal
O}(\qq)^{\pm 1}$ 成立，那么对 $\sum_i|n_i|$ 作归纳即可得到它的一般情形.

对于 ${\cal N}={\cal O}(\qq)^{-1}$，我们有 ${\cal L}\otimes {\cal
N}={\cal L}(-\qq)$，其截面是 ${\cal L}$ 在 $\qq$ 处消失的截面. 我们可以利用如下层的正合列\index{exact-sequence@正合列}来比较 ${\cal L}$ 与 ${\cal L}(-\qq)$ 的上同调：
\begin{equation}\label{doglo}
0\to {\cal L}(-\qq)\to {\cal L}\underarrow{r} {\cal L}|_\qq\to 0
\end{equation}
我们的记号略有些不正式. 我们用 ${\cal L}(-\qq)$ 和 ${\cal L}$ 表示所指的线丛，也表示相应的截面层\index{sheaf@截面层}；而 ${\cal L}|_\qq$ 既表示 ${\cal L}$ 在 $\qq$ 处的纤维，也表示 $\qq$ 处相应的摩天楼层\index{skyscraper-sheaf@摩天楼层}. 式（\ref{doglo}）中的映射 $r$ 通过在点 $\qq$ 处对 ${\cal L}$ 的截面取值来定义. 由式（\ref{doglo}），我们得到上同调群的长正合列
\begin{align}\label{zonglo}
 0\to & \,H^0(D,{\cal L}(-\qq))\to H^0(D,{\cal L})\to {\cal L}|_\qq \cr
 \to&\,
 H^1(D,{\cal L}(-\qq))\to H^1(D,{\cal L})\to 0.\end{align}
一般而言，向量空间的长正合列
\begin{equation}\label{onglo}
0\to A_0\underarrow{d_0}
A_1\underarrow{d_1}A_2\underarrow{d_2}\dots
\underarrow{d_{n-1}}A_n\to 0  \end{equation} 确定一个同构
\begin{equation}\label{ponglo}
\otimes_{i=0}^n \left(\det\,
A_i\right)^{(-1)^i}\cong\Bbb{C}.\end{equation} 在当前情形中，由式（\ref{zonglo}），我们得到同构
\begin{equation}\label{monglo}
\det\,H^*(D,{\cal L}(-\qq))\cong \det\,H^*(D,{\cal L})\otimes {\cal
L}|_\qq.\end{equation} 这等价于
\begin{equation}\label{kongo}
\langle{\cal L},{\cal O}(-\qq)\rangle \cong{\cal L}|_\qq^{-1}\otimes
{\det\,H^*(D,{\cal O}(-\qq))\over \det\,H^*(D,{\cal
O})}.\end{equation} 取 ${\cal L}={\cal O}$ 的特殊情形可知
\begin{equation}\label{ubongo}{\det\,H^*(D,{\cal O}(-\qq))\over \det\,H^*(D,{\cal O})}
\cong{\Bbb{C}},
\end{equation}
因此实际上
\begin{equation}\label{rongo} \langle{\cal L},{\cal O}(-\qq)\rangle\cong {\cal L}|_\qq^{-1}.
\end{equation}

如果在式（\ref{monglo}）中把 ${\cal L}$ 替换为 ${\cal M}$ 或 ${\cal L}\otimes {\cal M}$，则可知
\begin{align}\label{moonglo}
\det\,H^*(D,{\cal M}(-\qq))&\cong \det\,H^*(D,{\cal M})\otimes {\cal
M}|_\qq\\ \nonumber \det\,H^*(D,{\cal L}\otimes {\cal M}(-\qq))&\cong
\det\,H^*(D,{\cal L}\otimes{\cal M})\otimes ({\cal L}\otimes {\cal
M})|_\qq.
\end{align}

现在，如果写出 ${\cal N}={\cal O}(-\qq)$ 时 $\langle {\cal L},{\cal M}\otimes {\cal N}\rangle$ 的定义，利用式（\ref{moonglo}）把其中所有 ${\cal M}(-\qq)$ 消去而改用 ${\cal M}$，再通过式（\ref{rongo}）计算 ${\cal L}|_\qq^{-1}$，就得到
\begin{equation}\label{onear}\langle{\cal L},{\cal M}(-\qq)\rangle\cong\langle{\cal L},
{\cal M}\rangle\otimes \langle {\cal L},{\cal O}(-\qq)\rangle,\end{equation}
这就是 ${\cal N}={\cal O}(-\qq)$ 时的式（\ref{duglog}）.

如果在式（\ref{onear}）中令 ${\cal M}={\cal O}(\qq)$，则可知
\begin{equation}\label{gonear}\langle{\cal L},{\cal O}(\qq)\rangle\cong \langle{\cal L},
{\cal O}(-\qq)\rangle^{-1}.
\end{equation}
在式（\ref{onear}）中把 ${\cal M}$ 替换为 ${\cal M}(\qq)$，并利用式（\ref{gonear}），便得到 ${\cal N}={\cal O}(\qq)$ 时的式（\ref{duglog}）. 因此，由归纳法，这一结果一般地成立. 不过，我们应当指出，所给的解释稍显朴素，因为我们尚未证明：无论选择怎样把 $\N$ 表示为基本因子 $\O(\qq_i)^{\pm 1}$ 的张量积，都会得到同一个同构. 完整的论述可见 \cite{deligne}.

最后，我们注意到，这一结果可以写得更加对称. 对于任意线丛 ${\cal L}\to D$，将 $\det\,H^*(D,{\cal L})$ 简记为 $[{\cal L}]$. 为进一步简化记号，再略去张量积的符号. 那么，式（\ref{duglog}）就相当于说，对于任意三个线丛 ${\cal L}$、${\cal M}$、${\cal N}$，都有典范同构
\begin{equation}\label{hurglo} {[{\cal L}\,{\cal M}\,{\cal N}]\,[{\cal L}]\,[{\cal M}]\,
[{\cal N}]
\over [{\cal L}\,{\cal M}]\,[{\cal M}\,{\cal N}]\,[{\cal L}\,{\cal N}]\,[{\cal O}]}
\cong\Bbb{C}.
\end{equation}
这称为立方体定理. 这一名称反映了这样一个事实：左边的八个因子可以方便地排列在一个立方体的各个顶点上.

\subsection{解释}\label{interpretation}

现在，我们可以在一般情形下计算 $\langle{\cal L},{\cal M}\rangle$. 假设 ${\cal M}=\otimes_{i=1}^k{\cal O}(\qq_i)^{n_i}$，换言之，${\cal M} ={\cal O}({\DD})$，其中 ${\DD}=\sum_in_i\qq_i$ 是次数为 $d=\sum_in_i$ 的除子. 根据式（\ref{duglog}），并利用特殊情形（\ref{rongo}）和（\ref{gonear}），对 $\sum_i|n_i|$ 作归纳可得
\begin{equation}\label{trifong}
\langle{\cal L},{\cal M}\rangle=\otimes_{i=1}^k {\cal L}|_{\qq_i}^{n_i},\end{equation}
其中，与通常一样，${\cal L}|_{\qq_i}$ 是 ${\cal L}$ 在点 $\qq_i$ 处的纤维.

假设 ${\cal L}$ 的次数为 $c$. 我们希望固定 ${\cal M}$，让 ${\cal L}$ 变化. ${\cal L}$ 确定一个点 $x\in {\rm Pic}_c(D)$，为了明确表示这一点，我们将把 ${\cal L}$ 写为 ${\cal L}_x$. 随着 $x$ 变化，我们希望把 $\langle{\cal L},{\cal M}\rangle$ 解释为 ${\rm Pic}_c(D)$ 上一个线丛的纤维. 这里有一个关键的微妙之处，在考虑线丛的模空间\index{moduli-space@模空间}时常常会遇到它. 也就是说，一个点 $x\in {\rm Pic}_c(D)$ 只在相差同构的意义下确定一个线丛 ${\cal L}_x\to D$. ${\cal L}_x$ 可以替换为 ${\cal L}_x\otimes {\cal R}$，其中 ${\cal R}$ 是一个固定的一维向量空间. 从下面的意义上说，这就是改变 ${\cal L}_x$ 的一般自由度. 如果 ${\cal U}$ 和 ${\cal U}'$ 是黎曼曲面 $D$ 上的同构线丛，那么二者之间没有典范同构（因为它们各自都有一个 $\Bbb{C}^*$ 自同构群，该群中的元素可以与 ${\cal U}$ 和 ${\cal U}'$ 之间任意给定的同构复合）. 然而，存在一个典范同构 ${\cal U}'\cong {\cal U}\otimes {\cal R}$，其中 ${\cal R}$ 是一维向量空间 $H^0(D,{\cal U}'\otimes {\cal
U}^{-1})$. （构造这一同构，只需注意到：根据 $H^0(D,{\cal U}'\otimes {\cal U}^{-1})$ 的定义，${\cal R}$ 中的一个向量给出从 ${\cal U}$ 到 ${\cal U}'$ 的一个全纯映射.）因此，一旦给定 ${\cal L}_x$ 的同构类，它在相差 ${\cal L}_x\to {\cal L}_x \otimes {\cal R}$ 的意义下就是唯一的.

由式（\ref{trifong}），我们有
\begin{equation}\label{iftong}
\langle{\cal L}_x\otimes{\cal R},{\cal M}\rangle =\langle{\cal L}_x,{\cal M}\rangle
\otimes{\cal R}^d.
\end{equation}
因此，当且仅当 $d=0$ 时，$\langle{\cal L}_x,{\cal M}\rangle$ 才由 ${\cal L}_x$ 的同构类典范地确定. 一个相关的说法是：${\cal L}_x$ 上作用为乘以复数 $\lambda\in\Bbb{C}^*$ 的自同构，在 $\langle{\cal L}_x,{\cal M}\rangle$ 上作用为乘以 $\lambda^d$. 因此，当且仅当 ${\cal M}$ 的次数为零时，${\cal L}_x$ 的自同构才在 $\langle{\cal L}_x,{\cal M}\rangle$ 上平凡地作用.

所以，当 $d=0$ 时，对于每个 ${\cal M}=\otimes_i \O(\qq_i)^{n_i}$，我们确实得到 ${\rm Pic}_c(D)$ 上的一个线丛 ${\cal N}$. 这个线丛在 $x\in {\rm Pic}_c(D)$ 处的纤维是
\begin{equation}\label{zudco}
{\cal N}|_x=\otimes_{i=1}^k {\cal L}_x|_{\qq_i}^{n_i}.\end{equation}
这是一个初等的定义，即使事先不讨论行列式线丛，我们也可以作出它. 从上述讨论中我们了解到，在相差自然同构的意义下，以这种方式定义的线丛 $\cal N$ 仅依赖于线丛 $\M=\O(\DD)$ 的同构类，而不依赖于我们用来表示这一线丛的具体除子 $\DD=\sum_i n_i\qq_i$.

如果 ${\cal M}$ 的次数 $d$ 不为零，会发生什么？要作出类似的定义，我们必须先选取 $D\times \Pic_c(D)$ 上的一个万有线丛. 回顾第 \ref{eleig} 节中的这一概念：这样的万有线丛是一个线丛 $\W\to D\otimes \Pic_c(D)$，对于任意 $x\in \Pic_c(D)$，它在 $D\times x$ 上的限制具有与 $x$ 对应的同构类型. 这样的万有线丛是存在的，但其唯一性只到如下变换为止：
\begin{equation}\label{doffo}\W\to \W\otimes \pi_2^*(\RR),\end{equation}
其中 $\RR\to \Pic_c(D)$ 是某个线丛，而 $\pi_2:D\times\Pic_c(D)\to \Pic_c(D)$ 是投影.

式（\ref{zudco}）可以用万有丛\index{universal-bundle@万有丛} $\W$ 漂亮地改写. 给定 $\M=\otimes \O(\qq_i)^{n_i}$，$\Pic_c(D)$ 上对应的线丛为
\begin{equation}\label{offo} \N=\otimes_i \W|^{n_i}_{\qq_i\times \Pic_c(D)}. \end{equation}
要由此公式恢复式（\ref{zudco}），只需把 $\W|_{\qq_i\times x}$ 与 $\L_x|_{\qq_i}$ 等同起来. 对于 $\M$ 具有非零次数 $d$ 的情形，这就是我们能够对式（\ref{zudco}）作出的最佳推广. 然而，如果对于某个线丛 $\RR\to \Pic_c(D)$，把 $\W$ 变换为 $\W\otimes \RR$，则 $\N$ 相应变换为
\begin{equation}\label{woffo}\N\to \N\otimes \RR^d. \end{equation}
这种对万有丛选择的依赖意味着，$\N$ 实际上并未自然地定义为 $\Pic_c(D)$ 上的线丛，而应当更自然地理解为一个由某个 gerbe 扭曲的扭曲线丛. （这是一种同义反复：我们可以在 $\Pic_c(D)$ 上定义一个 gerbe，使得 $D\times \Pic_c(D)$ 上任意选取的万有线丛都自然地给出它的平凡化\index{trivialization@平凡化}；于是式（\ref{woffo}）就意味着，$\N\to \Pic_c(D)$ 最好理解为由该 gerbe 的 $d^{th}$ 次幂扭曲的扭曲线丛. 见第 \ref{gerbes} 节.）希望到第 \ref{compheck} 节时，这一表述的动机会变得清楚\transnote{局部地看，若把万有线丛换为 $W\otimes R$（$R$ 为 $\Pic_c(D)$ 上的线丛），式（\ref{woffo}）表明 $\N$ 的转移函数只在一个 $d$ 次上闭链（2-cocycle）意义下良定义；这正是“被 gerbe 的 $d$ 次幂扭曲的扭曲线丛”。$d$ 是 $\M$ 的次数，即总电荷，所以 §7 计算的赫克“本征值”只能是这种扭曲意义下的值；要得到真正的数，须取 $d=0$ 或先固定一个平凡化。[来源：原文 §6.2–§6.3；[KW]（arXiv:hep-th/0604151）§7；nLab "gerbe"，\url{https://ncatlab.org/nlab/show/gerbe}（2026-11-19 访问）。]}.





到目前为止，我们的讨论是不对称的. 我们考虑了符号 $\langle{\cal L},{\cal M}\rangle,$ 其中 ${\cal L}$ 和 ${\cal M}$ 的次数分别为 $c$ 和 $d$. 我们固定 ${\cal M}$ 而让 ${\cal L}$ 变化，并发现，在这种情况下，$\langle{\cal L},{\cal M}\rangle$ 随之变化，构成 ${\rm Pic}_c(D)$ 上一个扭曲次数为 $d$ 的扭曲线丛. 显然，由对称性可知，如果固定 ${\cal L}$ 而让 ${\cal M}$ 变化，我们将得到 ${\rm Pic}_d(D)$ 上一个扭曲次数为 $c$ 的扭曲线丛. 因此，当我们让 ${\cal L}$ 和 ${\cal M}$ 都变化时，就得到 ${\rm Pic}_c(D)\times {\rm Pic}_d(D)$ 上一个双次数为 $(d,c)$ 的扭曲线丛. 在对 gerbe 作一段插叙之后，我们将用对偶性的语言重新表述这一论断.

\subsection{Gerbe}\label{gerbes}

“Gerbe”已在本文的几处出现过. 我们将对这一概念作最低限度的解释，只说明可能有助于理解本文的最简单要点（尤其是为第 \ref{ugduality} 节提供背景）.

对于我们的目的，“gerbe”与群 $\C^*$ 或其子群 $U(1)$、$\Z_n$ 相联系. 空间 $X$ 上的 gerbe $\eG$ 在局部是平凡的，但在整体上可能并不平凡. （我们感兴趣的大多数 gerbe 在整体上都是平凡的，但并非典范地平凡.）一个 gerbe 在开集 $U\subset X$ 上的两个局部平凡化相差于张量乘以某个线丛 $\L\to U$. 我们不以抽象的方式解释这句话，但在例子中，它的含义会变得清楚. 对于 $U(1)$ gerbe，$\L$ 总带有一个厄米度量（因此它的结构群约化为 $U(1)$）. 对于 $\Z_n$ gerbe，$\L$ 总带有一个同构 $\L^n\cong \O$. 如果 gerbe 带有联络，那么与平凡化的局部改变相联系的线丛 $\L$ 也将带有联络.

下面是本文中已经出现的例子. 如果 $C$ 是一个黎曼曲面，那么 $C$ 上存在脚注 \ref{brief} 中提到的一个 $\Z_2$ gerbe，其平凡化是 $C$ 上的自旋结构\index{spin-structure@自旋结构}. 描述什么是自旋结构颇为微妙，但两个自旋结构相差于用一个二阶线丛 $\L$ 作扭曲，因此这是一个 $\Z_2$ gerbe. 整体自旋结构是存在的，所以这个 gerbe 是平凡的，但并非典范地平凡\transnote{$C$ 上的自旋结构可视为典范丛 $K$ 的平方根（theta characteristic）；任意两个相差一个二阶线丛，这正对应 $\Z_2$ gerbe 的“带”。对一族狄拉克算符（时空维数 $8k+2$ 时），行列式线丛有典范平方根（Pfaffian 线丛），但把它识别为某个线丛需要先选自旋结构，故平凡化非典范；这与脚注 \ref{brief} 中“希钦纤维的等同依赖 $K^{1/2}$ 的选择”是同一现象。[来源：原文 §6.3 与 §6.5 脚注；nLab "determinant line bundle" 之 "Pfaffian line bundle" 与 "square roots" 表（canonical bundle 的平方根 = theta characteristic = spin structure）（2026-11-19 访问）。]}. 作为第二个例子，考虑对某个 $c\in \Z$，$C\times \Pic_c(C)$ 上的万有线丛. 任意两个这样的万有线丛相差于张量乘以一个线丛 $\mathcal R\to \Pic_c(C)$；这告诉我们，$\Pic_c(C)$ 上存在一个 gerbe，$C\times \Pic_c(C)$ 上任意选取的万有线丛都使它平凡化. 这样的万有线丛是存在的，所以所讨论的 gerbe 是平凡的，但并非典范地平凡. 为给出一个涉及非平凡 gerbe 的类似例子，令 $^L\neg G$ 为具有非平凡中心 $\Z_n$ 的单非阿贝尔群. 那么，在 $^L\neg G$ 丛的模空间 $\M(^L\neg G,C)$ 上存在一个 gerbe，\footnote{如果将 $\M(^L\neg G,C)$ 替换为相应的希格斯丛\index{higgs-bundle@希格斯丛}模空间 $\MH(^L\neg G,C)$，我们接下来要作出的所有陈述仍然适用.}在开集 $U\subset \MH(^L\neg G,C)$ 上，选择 $C\times U$ 上的一个万有 $^L\neg G$-丛便使它平凡化. 这样的万有丛在局部存在，而且任意两种选择相差于用一个 $n$ 阶线丛 $\L\to U$ 作扭曲. 因此，在这一情形中寻找万有丛的问题，在 $\MH(^L\neg G,C)$ 上定义了一个 $\Z_n$ gerbe. 由于整体万有丛不存在，这个 gerbe 在拓扑上是非平凡的. 关于这个例子的更多讨论，见 \cite{KW} 的第 7 节. 这是与本文相关的唯一非平凡 gerbe\transnote{此 gerbe 的“带”（band）是中心 $\Z_n$：$^L\neg G$ 的万有丛在整体上不存在，因为中心在表示纤维上非平凡地作用；任意两个局部万有丛相差一个 $n$ 阶线丛的扭曲。故其障碍类落在 $H^2(\M(^L\neg G,C),\Z_n)$，非平凡性即“离散扭转”（discrete torsion）；[KW] 第 7 节正是用这一点研究 $^L\neg G$-丛的扭曲向量丛。[来源：[KW]（arXiv:hep-th/0604151）§7；原文 §6.3；nLab "gerbe" 之 "Banded gerbes" 节（2026-11-19 访问）。]}.

如果空间 $X$ 上的 gerbe $\eG$ 在整体上是平凡的，那么它的一个平凡化就称为扭曲线丛. 这一术语的动机如下. 两个平凡化之差是一个线丛 $\L\to X$，但任意单个平凡化并不是这样的线丛，因此我们将它称为由 $\eG$ 扭曲的扭曲线丛.\footnote{如果 $\eG$ 是一个非平凡 gerbe，那么根据定义，$\eG$-扭曲线丛不存在. 然而，存在 $\eG$-扭曲向量丛\index{vector-bundle@向量丛}的概念，而非平凡 gerbe 可能容许扭曲向量丛. 例如，给定 $^L\neg G$ 的一个表示 $^L\neg R$，且 $^L\neg G$ 的中心在该表示上非平凡地作用，那么万有丛 $\EE_{^L\neg R}\to C\times\M(^L\neg G,C)$ 作为向量丛并不存在，但作为 $\eG$-扭曲向量丛确实存在，其中 $\eG$ 是通过选取万有丛而在局部平凡化的 gerbe. 这一陈述是一种花哨的同义反复. 初等的解释见 \cite{KW} 的第 7 节.}

现在来讨论 $U(1)$ gerbe 上的联络. 如果 $U(1)$ gerbe $\eG$ 已经平凡化，那么联络可以简单地表示为一个二形式 $B$. （这类似于这样一个事实：如果线丛 $\L\to X$ 已经平凡化，那么 $\L$ 上的联络可以表示为一个一形式 $A$. 更一般地，非平凡 gerbe 上也可以定义联络，但我们不需要这一概念.）改变平凡化，是通过用一个带有联络 $A$ 的线丛 $\L\to X$ 作扭曲来实现的. 令 $F=dA$ 为相应的曲率. 在由 $\L$ 给出的平凡化改变下，二形式 $B$ 按 $B\to B+F$ 平移. 这里，$B+F$ 是 $\eG$ 相对于新平凡化的联络形式. 所谓平坦的 $\eG$-扭曲线丛（或等价地，$\eG$ 的平坦平凡化），是指使新的联络形式 $B+F$ 恒为零的平凡化. 为使平坦的 $\eG$-扭曲线丛存在，三形式曲率 $H=dB$（它在 $B\to B+F$ 下不变）必须消失，此时我们称 $\eG$ 为平坦 gerbe；并且二形式 $B/2\pi$ 的周期必须为整数\transnote{这里的 $B$ 即物理中的 $B$ 场（2-形式势）：改变平凡化时 $B\to B+F$ 是“大规范变换”，$H=dB$ 是规范不变的三形式场强；$B/2\pi$ 的周期取整为通量量子化条件。平坦 gerbe（$H=0$）上的平坦扭曲线丛，就是低能有效理论中“平坦 $B$ 场”的几何对应物。[来源：原文 §6.3；[KW]（arXiv:hep-th/0604151）§7；nLab "gerbe"（2026-11-19 访问）及其子条目 "gerbe (in differential geometry)"。]}. 如果存在平坦的 $\eG$-扭曲线丛，那么选取其中一个的方式可能并不典范，但任意两个都相差于用一个普通的平坦线丛\index{flat-line-bundle@平坦线丛}作扭曲.




\subsection{对偶性}
\label{ugduality}

在第 \ref{dualtori} 节中，我们引入了复环面 $\cF$ 与 $\cF'$ 之间的对偶性概念. 这样的对偶性意味着，$\cF$ 参数化 $\cF'$ 上的平坦线丛，反之亦然. 表述这种对偶性最有用的方式，是展示一个庞加莱线丛\index{poincare-line-bundle@庞加莱线丛}，即线丛 ${\cal T}\to \cF\times \cF'$，其主要性质是：对于 $f\in \cF$，它在 $f\times {}\cF'$ 上的限制，就是由 $f$ 标记的 $\cF'$ 上的线丛；交换 $\cF$ 与 $\cF'$ 的角色后也同样成立. 因此，随着 $f$ 变化，在 $f\times \cF'$ 上的限制给出 $\cF'$ 上平坦线丛的一个万有族，交换两个因子后也同样成立.

我们可以将这一概念推广为扭曲对偶性. 假设两个环面 $\cF$ 和 $\cF'$ 分别配备平坦的 $U(1)$ gerbe ${\eG}$ 和 ${\eG}'$，并假设这些平坦 gerbe 是平凡的（作为平坦 gerbe 平凡，而不仅仅是在拓扑上平凡）. 因此，通过选取一个扭曲线丛，${\eG}$ 或 ${\eG}'$ 都可以平凡化，但这种平凡化并不典范；而且，我们可以把这个扭曲线丛选为平坦的. 我们所谓 $\cF$ 与 $\cF'$ 之间类型为 $({\eG},{\eG}')$ 的扭曲对偶性，将 $\cF'$ 视为 ${\eG}$ 的平坦平凡化的模空间，并将 $\cF$ 视为 ${\eG}'$ 的平坦平凡化的模空间. 与普通情形一样，描述扭曲对偶性最自然的方式，是展示一个由 $\cF\times {}\cF'$ 上的 gerbe ${\eG}\otimes {}{\eG}'$ 扭曲的扭曲庞加莱线丛；随着 $f'$（或 $f$）变化，它在 $\cF\times{}f'$（或 $f\times {}\cF'$）上的限制，给出 ${\eG}$（或 ${\cal G}'$）的平坦平凡化的一个万有族.

我们断言，符号 $\langle{\cal L},{\cal M}\rangle$ 在 ${\rm Pic}_c(D)\times {\rm Pic}_d(D)$ 上给出一个双次数为 $(d,c)$ 的典范扭曲庞加莱线丛. 由于对应 ${\cal L},{\cal M}\to \langle{\cal L},{\cal M}\rangle$ 显然是典范的，而且以适当的方式扭曲，我们只需证明它确实是一个对偶性. 一般情形可以非典范地映到 $c=d=0$ 的情形；只需选取基点（即选取次数为 $c$ 或 $d$ 的特定平坦线丛，将其映到 ${\rm Jac}(D)$ 的原点）以及 gerbe 的平凡化，就可以把 ${\rm Pic}_c(D)$ 和 ${\rm Pic}_d(D)$ 映到 ${\rm Pic}_0(D)={\rm Jac}(D)$. 因此，要说明 $\langle{\cal L},{\cal M}\rangle$ 定义了次数为 $(d,c)$ 的典范扭曲对偶性，我们只需证明，在 $c=d=0$ 时，它确实定义了一个对偶性.

由于这是一个重要的事实，我们将用两种方式来解释它. 首先，在微分几何中，要描述黎曼曲面 $D$ 上的一族全纯线丛，可以考虑一个固定的光滑复线丛 ${\cal U}$，并在其上赋予一族酉联络. 如果把 ${\cal U}$ 上的联络记为 $U$，那么让 $U$ 变化所得到的一族复线丛上的行列式线丛，具有自然的厄米度量，以及曲率如下的自然酉联络 \cite{quillen}：
\begin{equation}\label{pliko}
{i\over 4\pi}\int_D \delta U\wedge \delta U.\end{equation} 现在回顾 $\langle{\cal L},{\cal M}\rangle$ 的定义：
\begin{equation}\label{zunko}
\langle{\cal L},{\cal M}\rangle={\det\,H^*(D,{\cal L})\otimes
\det\,H^*(D,{\cal M})\over \det\,H^*(D,{\cal L}\otimes {\cal
M})\otimes \det\,H^*(D,{\cal O})}.\end{equation} 我们把 ${\cal L}$ 和 ${\cal M}$ 视为固定的光滑线丛，其联络分别为 $A$ 和 $B$. 那么 ${\cal L}\otimes{\cal M}$ 就是带有联络 $A+B$ 的固定光滑线丛. ${\cal L}$、${\cal M}$ 和 ${\cal L}\otimes{\cal M}$ 的行列式线丛带有自然联络，其曲率分别通过在式（\ref{pliko}）中用 $A$、$B$ 或 $A+B$ 替代 $U$ 来得到. 然后，根据 $\langle{\cal L},{\cal M} \rangle$ 的定义，对这些曲率取适当的线性组合，就可知 $\langle{\cal L},{\cal M}\rangle$ 所给出的线丛带有一个自然联络，其曲率为
\begin{equation}\label{likko}
-{i\over 2\pi}\int_D\delta A\wedge \delta B.\end{equation}


到目前为止，我们得到的是 $\AA\times \AA'$ 上的一个线丛，其中两个因子分别是 ${\cal L}$ 和 ${\cal M}$ 上所有联络组成的空间. 我们希望将它下降为 ${\rm Jac}(D)\times{\rm Jac}(D)$ 上的一个线丛. ${\cal L}$ 和 ${\cal M}$ 的规范变换群作用在 $\AA\times \AA'$ 上，并且这一作用提升为线丛 $\langle{\cal L},{\cal M}\rangle$ 上的作用. 我们希望对规范群取商. 为此，必须要求 ${\cal L}$ 和 ${\cal M}$ 的次数为零. 否则，正如式（\ref{iftong}）之后解释的那样，在 $A$ 和 $B$ 上平凡作用的常数规范变换，将在 $\langle{\cal L},{\cal M}\rangle$ 上非平凡地作用，因而不能得到合理的商. 一旦将 ${\cal L}$ 和 ${\cal M}$ 限制为次数为零，就可以将联络 $A$ 和 $B$ 限制为平坦联络\index{flat-connection@平坦联络}；此时，如式（\ref{ivo}）之后所解释的，曲率形式（\ref{likko}）便下降到对规范变换群所取的商上. 这样，我们可以对作用在线丛 $\langle{\cal L},{\cal M}\rangle$ 及其联络上的规范变换取商，得到 ${\rm Jac}(D)\times {\rm Jac}(D)$ 上的一个线丛，其曲率仍可用相同的公式（\ref{likko}）表示. 正如讨论式（\ref{ivo}）时所指出的，${\rm Jac}(D)\times {\rm Jac}(D)$ 上具有这一曲率的酉线丛，就是一个庞加莱线丛\transnote{式（\ref{pliko}）是 Quillen 度量的曲率公式：对固定的光滑复线丛取一族酉联络（即一族 $\bar\partial$-算符），其行列式线丛带自然厄米度量与自然酉联络，曲率即式（\ref{pliko}）。$\langle{\cal L},{\cal M}\rangle$ 的曲率由式（\ref{zunko}）中四个行列式线丛的曲率按分式的 $\pm$ 取线性组合得到，即式（\ref{likko}）。当 ${\cal L},{\cal M}$ 限制为平坦联络、即 $\Jac(D)$ 上的点时，这个 2-形式是 $\Jac(D)\times\Jac(D)$ 的辛（Kähler）形式（相差归一化），故具此曲率的酉线丛必为庞加莱线丛。[来源：原文 §6.4；Quillen 1985（Funct. Anal. Appl. 19，31–34，本文文献 [41]）；2026-11-19 核对。]}.

另一种更接近代数几何的方法是，选取基点 $Q\in D$，并考虑嵌入 $\theta:D\to{\rm Jac}(D)$，它把 $R\in D$ 映到 ${\rm Jac}(D)$ 中与零次线丛 ${\cal O}(R)\otimes{\cal O}(Q)^{-1}$ 对应的点. 由于这一映射给出同构 $H_1(D,\Bbb{Z}) \cong H_1({\rm Jac}(D),\Bbb{Z})$，而平坦线丛由从 $H_1$ 到 $U(1)$ 的同态分类，因此，平坦线丛 $W\to{\rm Jac}(D)$ 与平坦线丛 ${\cal M}\to D$（或等价地，次数为零的全纯线丛 ${\cal M}\to D$）之间存在一一对应 $W\to \theta^*(W)$.

为建立对偶性，我们固定 ${\cal M}$，让 ${\cal L}$ 变化，并将对应 ${\cal L}\to \langle{\cal L},{\cal M}\rangle$ 视为在 ${\rm Jac}(D)$ 上定义了一个线丛 $W_{\cal M}$. 我们希望说明，在相差同构的意义下，每个平坦线丛 $W\to {\rm Jac}(D)$ 都恰好与一个 $W_{\cal M}$ 同构. 但根据上一段的说明，只需证明，对于某个唯一的 ${\cal M}$，有 $\theta^*(W)=\theta^*(W_{\cal M})$. 我们断言，对于任意零次线丛 ${\cal M}\to D$，都有
\begin{equation}\label{yormy}\theta^*(W_{\cal M})\cong {\cal M}.\end{equation}
因此，当且仅当取 ${\cal M}=\theta^*(W)$ 时，才有 $\theta^*(W)\cong \theta^*(W_{\cal M})$.

为证明断言（\ref{yormy}），我们必须理解 ${\cal L}={\cal O}(R)\otimes {\cal O}(Q)^{-1}$ 时的 $\langle{\cal L},{\cal M}\rangle$. 为此，只需利用式（\ref{trifong}），但这次交换 ${\cal L}$ 与 ${\cal M}$ 的角色，便知 $\langle{\cal O}(R)\otimes {\cal O}(Q^{-1}),{\cal M}\rangle \cong {\cal M}|_R\otimes {\cal M}|_Q^{-1}$. 我们希望固定 $Q$（因此 ${\cal M}|_Q^{-1}$ 是一个无关紧要的固定一维向量空间 ${\cal R}$），让 $R$ 变化. 在 $R$ 处的纤维为 ${\cal M}|_R\otimes {\cal M}|_Q^{-1}$ 的线丛，就是 ${\cal M}\otimes {\cal R}$. 换言之，它与 ${\cal M}$ 同构，这正是我们要证明的.

\subsection{再论希钦纤维的对偶性}
\label{dualap}

在谱覆盖\index{spectral-cover@谱覆盖} $\psi:D\to C$ 的语境下，$D$ 的雅可比簇\index{jacobian@雅可比簇}的自对偶性意味着，对于自对偶群 $U(N)$，希钦纤维化\index{hitchin-fibration@希钦纤维化}的纤维是自对偶的.\footnote{如第 \ref{unitaryexample} 节所解释的（见脚注 \ref{brief}），当 $N$ 为偶数时，将希钦纤维化的纤维与雅可比簇等同起来，依赖于对典范丛\index{canonical-bundle@典范丛} $K$ 的平方根的选择. 同样，我们下面使用的对 $SU(N)$ 纤维的描述，在 $N$ 为偶数时也依赖于自旋结构的选择.} 现在，利用谱覆盖 $\psi:D\to C$，重新考虑群 $SU(N)$ 和 $PSU(N)$ 的希钦纤维之间的对偶性.

我们已经解释过，对于 ${\cal L}\in {\rm Pic}_c(D)$、${\cal M}\in{\rm Pic}_d(D)$，符号 $\langle{\cal L},{\cal M}\rangle$ 给出 ${\rm Pic}_c(D)$ 与 ${\rm Pic}_d(D)$ 之间次数为 $(d,c)$ 的扭曲对偶性. 这就是 $G=U(N)$ 的希钦纤维的自对偶性，推广到了允许电通量和磁通量的情形\transnote{在谱覆盖 $\psi:D\to C$ 下，$U(N)$ 的希钦纤维（选定 $K^{1/2}$ 后）可等同于 $D$ 的雅可比簇（同构依赖 $K^{1/2}$ 的选择）；§6.4 的配对说明 $\Jac(D)$ 与其对偶阿贝尔簇可由庞加莱线丛等同，谱覆盖再把这一等同翻译成希钦纤维的对偶性。注意这里的纤维是阿贝尔簇（可带扭化），并非字面意义上的环面。[来源：原文 §6.5；Hitchin 1987（Proc. London Math. Soc. 55，59–126，本文文献 [13]）；2026-11-19 核对。]}. （这一解释将在第 \ref{abcase} 节中给出.）现在我们断言，如果将第一个变量限制在 $SU(N)$ 的希钦纤维中，那么第二个变量将自然地投影到 $PSU(N)$ 的情形，从而得到所需的 $SU(N)$ 与 $PSU(N)$ 之间的扭曲对偶性. 因此，要求 ${\cal L}$ 取值于 $\FF_{SU(N)}$；换言之，取 ${\cal L}\in{\rm Jac}(D)$，并要求 ${\rm Nm}({\cal L})$ 为平凡线丛. 正如马上将要证明的，如果 ${\cal M}_0$ 是 $C$ 上的一个零次线丛，那么 $\langle{\cal L},\psi^*({\cal M}_0)\rangle$（其中让 ${\cal L}$ 变化，而固定 ${\rm Nm}({\cal L})$ 和 ${\cal M}_0$）作为 $\FF_{SU(N)}$ 上的线丛是平凡的. 有了这一点，配对 $\langle~,~\rangle$ 的双线性就意味着，对于 ${\cal M}_0\in {\rm Jac}(C)$，有 $\langle{\cal L},{\cal M}\rangle\cong\langle{\cal L},{\cal M}\otimes \psi^*({\cal M}_0)\rangle$. 因而，在如此限制 ${\cal L}$ 后，我们可以认为 ${\cal M}$ 取值于 $\FF^{(d)}_{PSU(N)}={\rm Pic}_d(D)/{\rm Jac}(C)$. 于是，符号 $\langle{\cal L},{\cal M}\rangle$ 在 $\FF_{SU(N)}\times \FF^{(d)}_{PSU(N)}$ 上定义了一个次数为 $(d,0)$ 的扭曲线丛. 因此，$\FF_{SU(N)}$ 参数化 $\FF^{(d)}_{PSU(N)}$ 上的一族普通线丛，而 $\FF^{(d)}_{PSU(N)}$ 参数化 $\FF_{SU(N)}$ 上的一族扭曲次数为 $d$ 的扭曲线丛. 这就是预期的对偶性. 它之所以是对偶性而不仅仅是配对，依赖于这样一个事实：除了已经考虑到的对称性，符号 $\langle {\cal L},{\cal M}\rangle$ 对第二个变量没有额外的对称性；我们将在本节末尾简要讨论这一点.

我们可以进一步推广，让 ${\cal L}$ 的次数为 $c$. 现在选取一个固定的、次数为 $c$ 的线丛 ${\cal L}_0\to C$，并考虑满足 ${\rm Nm}({\cal L})={\cal L}_0$ 的 ${\cal L}\in {\rm Pic}_{c}(D)$. 我们断言，在用零次线丛 ${\cal M}_0\to C$ 的拉回\index{pullback@拉回}来扭曲 ${\cal M}$ 时，$\langle{\cal L},{\cal M}\rangle$ 仍然不变. 由此，符号 $\langle{\cal L},{\cal M}\rangle$ 给出 $\hat \FF^{(c)}_{PSU(N)}$ 与 $\FF^{(d)}_{PSU(N)}$ 之间所预期的次数为 $(d,c)$ 的对偶性.

在这些陈述中，整数 $c$ 和 $d$ 对应于 $PSU(N)$ 或 $SU(N)$ 规范理论\index{gauge-theory@规范理论}中可能出现的离散电荷\index{electric-charge@电荷}和磁荷\index{magnetic-charge@磁荷}.\footnote{读者可以先阅读第 \ref{compheck} 节中对 $G=U(1)$ 情形下同一问题更详细的解释，再回头阅读本段.} 回顾一下，对于这些群，可能存在取值于 $\Z_N$ 的离散电荷和磁荷. $\Pic_c(D)$ 参数化离散磁荷取值为 $c$ 的希格斯丛，而 $\Pic_c(D)$ 上由万有丛的 $d^{th}$ 次幂扭曲的扭曲线丛，可以作为一个离散电荷为 $d$ 的膜的陈-帕顿丛\index{chan-paton-bundle@陈-帕顿丛}. 在 $PSU(N)$ 与 $SU(N)$ 的对偶性中，$c$ 和 $d$ 的角色相互交换，正如所预期的那样.

还需要证明，如果固定 ${\rm Nm}({\cal L})$，那么随着 ${\cal L}$ 变化，$\langle{\cal L},\psi^*({\cal M}_0)\rangle$ 保持不变. 为此，需要知道几个事实. 首先，对于黎曼曲面之间的映射 $\psi:D\to C$，直接有 $H^*(D,{\cal L})=H^*(C,\psi_*{\cal L})$，因此
\begin{equation}\label{luto} \det\,H^*(D,{\cal
L})=\det\,H^*(C,\psi_*{\cal L}).\end{equation} 当然，右边就是 $\det\, H^*(C,E)$，其中 $E=\psi_*({\cal L})$. 此外，对于黎曼曲面 $C$ 上的任意向量丛 $E$，有
\begin{equation}\label{unofto}\det\,H^*(C,E)=\det\,H^*(C,\det\,E).\end{equation}
这可以对 $E$ 的秩作归纳来证明\transnote{式（\ref{unofto}）说 $\det H^*(C,E)$ 作为 $E$ 的函数只通过 $\det E$ 分解：右端秩为 1，只含 $\det E$ 的信息，而左端本可依赖 $E$ 的高阶陈类，故这一恒等式并不显然。证明用线子丛 ${\cal M}\subset E$、正合列（\ref{gonno}）与相应的长正合列，把秩 $r$ 的情形化到秩 $r-1$；§6.5 正是用它证明：固定 ${\rm Nm}({\cal L})$ 时 $\langle{\cal L},\psi^*{\cal M}_0\rangle$ 不随 ${\cal L}$ 变化。[来源：原文 §6.5；2026-11-19 核对。]}. 任意丛 $E\to C$ 都有一个全纯线子丛 ${\cal M}$，因而出现在如下正合列中：\begin{equation}\label{gonno} 0\to
{\cal M}\to E\to U\to 0.\end{equation} 由此，我们得到同构 $\det\, E\cong\det {\cal M}\otimes {\rm det}\,U$（与式（\ref{onglo}）的讨论相同）. 另外，我们可以从式（\ref{gonno}）导出一个上同调群的长正合列；如式（\ref{monglo}）的推导一样，它给出同构 $\det\,H^*(C,E)\cong \det\,H^*(C,{\cal M})\otimes\det\,H^*(C,U)$. 将这些结合起来，对 $E$ 的秩作归纳便得到式（\ref{unofto}）.

现在，对于线丛 ${\cal L}\to D$、${\cal M}_0\to C$，有 $\det\,H^*(D,{\cal L}\otimes \psi^*({\cal M}_0))=\det\,H^*(C,\psi_*({\cal L}\otimes \psi^*({\cal M}_0)))=\det\,H^*(C,E\otimes {\cal M}_0)=\det\,H^*(C,\det (E\otimes {\cal M}_0))$. 同样，$\det\,H^*(D,{\cal L})=\det\,H^*(C,\det E)$. 如果让 ${\cal L}$ 变化而固定 ${\rm Nm}({\cal L})$，那么 $\det E$ 也保持固定，因而最终 $\det\,H^*(D,{\cal L}\otimes {\cal M})\otimes\det\,H^*(D,{\cal L})^{-1}$ 也保持固定. 而在 $\langle{\cal L},\psi^*({\cal M}_0)\rangle$ 的定义中，这正是依赖于 ${\cal L}$ 的因子. 因此，当 ${\cal L}$ 变化而 ${\rm Nm}({\cal L})$ 和 ${\cal M}_0$ 保持固定时，$\langle{\cal L},\psi^*({\cal M}_0)\rangle$ 也保持固定，正如所断言的那样.

最后，来证明：如果对于零次的 ${\cal M}$，$\langle{\cal L},{\cal M}\rangle$ 对 ${\cal L}\in\FF_{SU(N)}$ 是平凡的，那么 ${\cal M}$ 就是从 $C$ 拉回的线丛. 我们在第 \ref{character} 节中证明，如果 ${\rm Nm}({\cal L})={\cal O}$，那么 ${\cal L}$ 是若干线丛 ${\cal O}(\qq)\otimes {\cal O}(\qq')^{-1}$ 的张量积，其中 $\psi(\qq)=\psi(\qq')$. 根据配对的双线性，$\langle{\cal L},{\cal M}\rangle$ 对这样的张量积 ${\cal L}$ 是平凡的，当且仅当它对单个因子是平凡的. 因此，令 ${\cal L}={\cal O}(\qq)\otimes {\cal O}(\qq')^{-1}$，得到 $\langle{\cal L},{\cal M}\rangle={\cal M}(\qq)\otimes {\cal M}(\qq')^{-1}$. 所以，随着 ${\cal L}$ 变化，$\langle{\cal L},{\cal M}\rangle$ 的一个平凡化，给出在 $\qq$ 和 $\qq'$ 位于 $C$ 中同一点之上时，将 ${\cal M}(\qq)$ 与 ${\cal M}(\qq')$ 等同的方式. 这种等同的存在，意味着 ${\cal M}$ 是从 $C$ 拉回的.

我们在这里隐含地使用了这样一个事实：参数化满足 $\psi(\qq)=\psi(\qq')$ 的点对 $\qq,\qq'\in D$ 的簇是不可约的. 可以通过寻找一个具有足够大单值变换\index{monodromy@单值变换}的谱曲线\index{spectral-curve@谱曲线}局部模型来说明这一点，例如 $y^N+\epsilon yz+z=0$，其中 $z$ 是 $C$ 上的局部参数，而 $\epsilon$ 是复常数.


\section{计算赫克本征值}
\label{compheck}

最后，我们将利用已经掌握的知识，计算 't Hooft 算符作用在磁本征膜\index{magnetic-eigenbrane@磁本征膜}上时的“本征值”. 我们在第 \ref{abcase} 节先从阿贝尔情形开始. 这是显然的出发点，也必然导向最直接的计算；不过，由于一个初等事实，我们会遇到一些解释上的微妙问题：与单非阿贝尔群不同，群 $U(1)$ 的中心和基本群都具有无限阶\transnote{对 $U(1)$ 来说中心与基本群都是 $\Z$：电荷 $d\in\Z$ 无界，磁通量（即 $c_1$）也可取任意整数，故 't Hooft 算符能把 $c_1({\cal L})$ 平移任意整数，把膜移出 $\MH$。对单非阿贝尔群，中心有限（如 $SU(N)$ 的中心为 $\Z_N$），可平移的 $c_1$ 只能是 $N$ 的倍数，于是磁通量成为离散不变量 $\xi$（§2.6.6）；这正是阿贝尔情形必须额外保留 $U(1)$ 规范多重态的根本原因。[来源：原文 §7、§7.1.2、§2.6.6；2026-11-19 核对。]}.

阿贝尔规范理论\index{gauge-theory@规范理论}中的几何朗兰兹对偶性，通常通过归于 Deligne 的一种略有不同的论证来描述，例如可见 \cite{frenkel}. 我们采用的途径，将为第 \ref{actsun} 节要讨论的非阿贝尔推广提供一个有用的出发点.

\subsection{阿贝尔情形}
\label{abcase}
\subsubsection{计算}

对于 $G=U(1)$，黎曼曲面 $C$ 上的希格斯丛\index{higgs-bundle@希格斯丛}就是一对 $({\cal L},\varphi)$，其中 ${\cal L}$ 是复线丛\index{line-bundle@线丛}，而 $\varphi\in H^0(C,K)$ 是全纯微分. $\varphi$ 在分析中实际上几乎不起作用，原因在于：由于 $U(1)$ 是阿贝尔群，希钦方程\index{hitchin-equation@希钦方程}是线性的，规范场\index{gauge-field@规范场}与希格斯场\index{higgs-field@希格斯场}彼此解耦，$\varphi$ 在赫克变换\index{hecke-transformation@赫克变换}下不变，而 $S$-对偶性的有趣作用仅仅发生在规范场上.

然而，有一个重要之处在第 \ref{unitary} 节中也讨论过：希钦模空间\index{hitchin-moduli-space@希钦模空间} $\MH(U(1),C)$ 参数化的是线丛 $\L$ 的次数为零的希格斯对 $(\L,\varphi)$. 一般而言，'t Hooft 算符的作用不保持这一条件，因此需要作一些讨论.



我们在三维流形 $W=I\times C$ 上，以哈密顿框架考虑 $U(1)$ 规范理论，其中 $I$ 是一个区间. $W$ 右端的边界条件\index{boundary-condition@边界条件}由一个膜定义，我们把它选为电本征膜\index{electric-eigenbrane@电本征膜}或磁本征膜. 我们用插入在点 $y_i\times \pp_i\in I\times C$ 处、带有电荷\index{electric-charge@电荷} $d_i$ 的威尔逊算符 $W^{(d_i)}(y_i,\pp_i)$ 作用在电本征膜 ${\cal B}$ 上. 或者，对偶地，用 't Hooft 算符 $T^{(d_i)}(y_i,\pp_i)$ 作用在磁本征膜 $\hat{\cal B}$ 上. 第 \ref{thooftheckeop} 节的图 \ref{nuxto} 示意了相应的二维有效图像.

具体而言，我们选取一个电本征膜 ${\cal B}$，它是支撑在某点 $x\in \MH$ 上的零维膜，该点对应于希格斯丛 $({\cal L},\varphi)$. 然后考虑作用在 ${\cal B}$ 上的一组威尔逊算符的乘积：
\begin{equation}\label{pilkob}{\prod_{i=1}^n} \,W^{(d_i)}(y_i,\pp_i)\,{\cal B}.\end{equation}
如第 \ref{eleig} 节所解释的，零维膜 ${\cal B}$ 是电本征膜，并且有
\begin{equation}\label{nilkob}{\prod_{i=1}^n} \,W^{(d_i)}(y_i,\pp_i)\,{\cal B}=
{\cal B}\otimes\left(\otimes_{i=1}^n{\cal \L}|_{\pp_i}^{d_i}\right).\end{equation}
（稍后我们会将其改写为一种更接近第 \ref{eleig} 节中式（\ref{zefor}）的形式.）

这里一个重要的微妙之处，来自这样一个熟知的事实：点 $x\in \MH$ 只在相差同构的意义下确定相应的线丛 ${\cal L}$. 我们可以自由地将 ${\cal L}$ 与一个固定的一维向量空间 ${\cal R}$ 作张量积. 如果 $d=\sum_id_i$ 等于零，这不会带来问题；但一般而言，给定的威尔逊算符乘积作用在 ${\cal B}$ 上时的“本征值”，会被张量乘以 ${\cal R}^d$. 当然，这种不确定性与点 $\pp_i$ 的选择无关，我们在第 \ref{interpretation} 节中已经以另一种形式遇到过它. 与那里的讨论一样，最有启发性的做法是选取万有丛\index{universal-bundle@万有丛} $\W\to C\times {\mathrm{Jac}}(C)$，并用这一万有丛来表示式（\ref{nilkob}）：
\begin{equation}\label{wilkob} {\prod_{i=1}^n} \,W^{(d_i)}(y_i,\pp_i)\,{\cal B}=
{\cal B}\otimes\left(\otimes_{i=1}^n{\W}|_{\pp_i\times x}^{d_i}\right).\end{equation}

对于 $d\not=0$，这一结果确实依赖于万有丛 $\W\to C\times {\mathrm{Jac}}(C)$ 的选择. 如果把 $\W$ 变换为 $\W\otimes \RR$，其中 $\RR$ 是将第二个因子上的一个线丛拉回\index{pullback@拉回}到 $C\times \mathrm{Jac}(C)$ 所得的线丛，那么式（\ref{wilkob}）的右边就会被张量乘以 $\RR^d$.

其物理意义如下. 整数 $d$ 是我们用来作用的威尔逊算符乘积 ${\prod_{i=1}^n} \,W^{(d_i)}(y_i,\pp_i)$ 所携带的总电荷. 如果初始膜 $\B$ 是电中性的，那么这一威尔逊算符乘积作用后所得的膜就带有电荷 $d$.\footnote{在有效的 $1+1$ 维描述中，膜是弦端点处的边界条件. 说一个膜 $\B$ 带有电荷 $d$，只是意味着在这一边界条件下，电荷算符在弦的端点处得到一个 $d$ 的贡献.} 靶空间为 $\MH(U(1),C)$ 的 $\sigma$-模型只描述中性自由度；要描述带电的膜，不能仅仅使用这个 $\sigma$-模型，而需要在有效的二维描述中包含未破缺的 $U(1)$ 规范多重态. 第 \ref{elmag} 节将对此作进一步讨论.

一般而言，$\MH(U(1),C)$ 上的中性膜具有陈-帕顿丛\index{chan-paton-bundle@陈-帕顿丛} $\U$，它是普通的向量丛\index{vector-bundle@向量丛}或层. 电荷为 $d$ 的膜则具有扭曲向量丛或扭曲层形式的陈-帕顿丛 $\U$. 描述这种扭曲最简单的方式是：一旦选取万有丛 $\W=C\times \mathrm{Jac}(C)$，就可以将 $\U$ 描述为普通的向量丛或层，但 $\U$ 依赖于 $\W$ 的选择：将 $\W$ 变换为 $\W\otimes \RR$，会把 $\U$ 变换为 $\U\otimes \RR^d$. 显然，在如此说明带有净电荷的膜的陈-帕顿丛是何种对象之后，式（\ref{wilkob}）就与我们预期威尔逊算符的作用改变膜的电荷的方式一致.


电本征膜 ${\cal B}$ 的对偶是一个 ${\bf F}$ 型膜 $\hat{\cal B}$，我们预期它是一个磁本征膜. 它支撑在希钦纤维化\index{hitchin-fibration@希钦纤维化}的某个纤维 $\FF$ 上，该纤维所对应的希格斯场 $\tilde \varphi $ 是电侧原始希格斯场 $\varphi$ 的一个倍数.\footnote{这一倍数为 $\mathrm{Im}\,\tau$，其中 $\tau=\theta/2\pi + 4\pi i/e^2$ 是通常的规范耦合参数. 一般可将它通过尺度变换消去，它在下文中不起实质作用.} $\FF$ 是雅可比簇\index{jacobian@雅可比簇} ${\rm Jac}(C)$ 的一个副本. 膜 $\hat \B$ 由一个平坦陈-帕顿线丛\index{chan-paton-line-bundle@陈-帕顿线丛}刻画，我们马上就会讨论它. 我们希望计算一组 't Hooft 算符的乘积在 $\hat{\cal B}$ 上的作用：
\begin{equation}\label{doppy}  {\prod_{i=1}^n} \,T^{(d_i)}(y_i,\pp_i)\,\hat{\cal B}.
\end{equation}
这里同样有一个微妙之处：'t Hooft 算符的乘积将线丛的次数平移 $d=\sum_id_i$，因此如果 $d$ 非零，我们就被映到了 $\MH$ 之外. 与电侧情形一样，这意味着膜可以携带磁通量，而靶空间为 $\MH$ 的二维 $\sigma$-模型并未包含这种磁通量.

对于单或半单的非阿贝尔群 $G$，类似的磁通量概念是希格斯丛的一个离散拓扑不变量 $\xi$（在第 \ref{psun} 节中引入）. 因为它是离散不变量，所以可以由平坦丛携带，也可以由 $G$ 的希钦方程的一个解携带. 因此，对于单或半单的 $G$，$\MH(G,C)$ 具有由 $\xi$ 标记的连通分支. 但对于 $G=U(1)$，在靶空间为 $\MH(G,C)$ 的 $\sigma$-模型中描述磁通量并不方便；要描述它，必须在低能描述中保留未破缺的 $U(1)$ 规范对称性. 我们将在第 \ref{elmag} 节解释，对于 $G=U(1)$，什么样的膜会携带磁通量.

暂且搁置解释方面的问题，计算 't Hooft 算符在 $\hat{\cal B}$ 上的作用实际上是一个简单的练习. 根据第 \ref{dualap} 节中对对偶性的重新考察，支撑在 $\Jac(C)$ 中与线丛 $\L\to C$ 对应的点上的零维膜，其对偶应当是一个 $\bf F$ 型膜 $\hat \B$，其陈-帕顿丛为 $\langle\L,\M\rangle$. 这里，固定 $\L$，将此表达式看作由 $\M$ 参数化的一个 $\Jac(C)$ 副本上的线丛. 根据式（\ref{welmo}），用所示的 't Hooft 算符乘积作用在 $\hat{\cal B}$ 上所得的膜 $\prod_{i=1}^n\,T^{(d_i)}(y_i,\pp_i)\,\hat{\cal B}$ 的陈-帕顿丛，是 $\tilde \U=\Phi^*(\U)$，其中 $\Phi$ 是 $\Jac(C)$ 的自同构 $\M\to \M\otimes \O(\pp_i)^{d_i}$.

有 $\Phi^*(\langle\L,\M\rangle)=\langle\L,\M\otimes_i \O(\pp_i)^{d_i}\rangle$. 利用符号 $\langle~,~\rangle$ 的双线性，它等于 $\langle\L,\M\rangle\otimes \langle\L,\otimes_i\O(\pp_i)^{d_i}\rangle$. 换言之，
\begin{equation}\label{menz}\hat\U=\U\otimes\langle\L,\otimes _i\O(\pp_i)^{d_i}\rangle. \end{equation}
这证实膜 $\hat \B$ 是一个磁本征膜，其“本征值”为一维向量空间 $\langle\L,\otimes _i\O(\pp_i)^{d_i}\rangle=\otimes_i \L_{\pp_i}^{d_i}$，这里我们利用式（\ref{trifong}）计算了符号 $\langle~,~\rangle$. 因此，膜 $\hat {\cal B}$ 是磁本征膜，并具有从电侧公式（\ref{nilkob}）所预期的“本征值”：
\begin{equation}\label{rgo}
{\prod_{i=1}^n} \,T^{(d_i)}(y_i,\pp_i)\,\hat{\cal B}=\hat{\cal B}\otimes
\left(\otimes_{i=1}^n{\cal L}|_{\pp_i}^{d_i}\right).
\end{equation}

这就建立了对偶性. 't Hooft 算符在 $\hat{\cal B}$ 上的作用，与威尔逊算符在 ${\cal B}$ 上的作用相同. 现在只需解释，当 $d=\sum_id_i$ 非零时，这些计算意味着什么.

\subsubsection{电场与磁场}
\label{elmag}

在用 't Hooft 算符作用时，不可避免地会出现满足 $c_1({\cal L})\not= 0$ 的希格斯对 $({\cal L},\varphi)$. 但它们并不对应于 $\MH$ 中的点. 那么，它们如何进入规范理论？它们对应于我们尚未考虑过的一类新膜.

这一基本的微妙之处实际上已在第 \ref{prelims} 节中指出. 关于在黎曼曲面 $C$ 上紧致化\index{compactification@紧致化}的四维规范理论会在低能下约化为靶空间是 $\MH(G,C)$ 的二维 $\sigma$-模型的论证，假设在所关心的情形中，规范对称性完全破缺. 对于单非阿贝尔规范群 $G$，只要避开 $\MH$ 的奇点，情况确实如此（只剩下有限群，即 $G$ 的中心）；在这些奇点处，新的自由度会变得相关. 然而，$G=U(1)$ 的情形不同，因为常数规范变换平凡地作用，而每个伴随值场都具有一个连续的规范对称群. 对规范群为 $U(1)$ 的 $\N=4$ 超杨-米尔斯理论\index{super-yang-mills-theory@超杨-米尔斯理论}，无论怎样从四维紧致化到二维，$U(1)$ 规范群都保持未破缺，因此应当将其纳入低能描述.

因此，对于 $G=U(1)$，约化为低能 $\sigma$-模型的说法并不完全成立. 低能理论是一个靶空间为 $\MH$ 的 $\sigma$-模型与一个规范群为 $U(1)$ 的超对称\index{supersymmetry@超对称}规范理论的乘积. 我们不能完全在 $\sigma$-模型中描述膜，还必须考虑规范理论.

考虑三维流形 $W=I\times C$ 上的扩展博戈莫尔尼方程\index{extended-bogomolny-equation@扩展博戈莫尔尼方程}\footnote{等价地，考虑 \cite{KW} 中在 $t=1$ 时的超对称方程（3.29），这是研究 $K$ 型 $A$-模型的一个方便取值.}（\ref{woff}）和（\ref{old}）. 我们将它们看作描述与 $W$ 上量子化相关的、不依赖于时间的超对称场构型的方程. 给 $W$ 赋予度量 $dy^2 +d\Omega^2$，其中 $y=x^1$ 是 $I$ 上的欧几里得坐标，而 $d\Omega^2$ 是 $C$ 上与 $y$ 无关的凯勒度量. 对于规范群 $U(1)$，这些方程存在一类与迄今所考虑的解相当不同的解. 令 $\omega$ 为 $C$ 的凯勒形式的某个倍数，并将其归一化为满足 $\int_C\omega = 1$. 我们可以在次数为 $d$ 的线丛 ${\cal L}\to C$ 上选取曲率为 $F=2\pi  d\,\omega$ 的联络 $A$，来求解扩展博戈莫尔尼方程.\footnote{为避免出现看起来不自然的 $i$ 因子，这里取 $F$ 和 $\phi$ 为实值. 因此，线丛上酉联络 $D$ 的曲率为 $F=-iD^2$.} 然后，将 ${\cal L}$ 和 $A$ 拉回到 $W$，并取 $\phi_0=2\pi y\cdot  d$.\footnote{这里，$\phi_0$ 只是希格斯场 $\phi=\sum_{i=0}^3\phi_i dx^i$ 的“时间”分量.} 将其他所有场都取为零. 这样，对于 $G=U(1)$，我们就得到一个超对称构型，其中 $C$ 上的线丛可以具有任意次数\transnote{可验算其归一化：由 $\int_C\omega=1$ 与 $F=2\pi d\,\omega$ 得 $\int_C F=2\pi d$，与 $c_1({\cal L})=d$ 的 Chern–Weil 表示一致；又 $\phi_0=2\pi y\,d$ 使 $\partial_y\phi_0=2\pi d$，正是后文“$\phi_0$ 的法向导数等于 $2\pi d$”的边界条件。故 $d\not=0$ 时此解给出带磁荷 $d$ 的膜，而非 $\MH$ 上的零维膜——低能描述必须含 $U(1)$ 向量多重态。[来源：原文 §7.1.2；[KW]（arXiv:hep-th/0604151）式（3.29）取 $t=1$；2026-11-19 核对。]}.

回到式（\ref{doppy}）中 $d=\sum_id_i$ 非零时会发生什么的问题，答案是：在这种情况下，在所有 't Hooft 算符的左侧，线丛 ${\cal L}$ 具有非零次数，并将基于刚才描述的解产生一个超对称构型. 因此，总 $d$ 非零的一组 't Hooft 算符的乘积作用在 $\hat {\cal B}$ 上时，会产生一个不同于迄今所考虑的膜. 这是一个超对称膜，其边界条件要求 $c_1({\cal L})=d$，并要求 $\phi_0$ 的法向导数等于 $2\pi  d$. 我们可以称它为磁荷\index{magnetic-charge@磁荷}为 $d$ 的膜.

对于半单非阿贝尔群 $G$，只要避开 $\MH$ 的奇点，这一构造就没有十分相近的对应，因为扩展博戈莫尔尼方程会迫使 $\phi_0$ 与所有其他场对易，并生成解的一个对称性. 要在单非阿贝尔群 $G$ 的情形下得到与 $U(1)$ 中发现的结果类似的东西，可以让 $C$ 的亏格为零或一；此时，一个一般的稳定希格斯丛 $(E,\varphi)$ 可以保留一个连续的未破缺规范对称群. 对于亏格为 0 的情形，$E$ 平凡且 $\varphi=0$ 的半稳定\index{semistability@半稳定}希格斯丛是一般的（作为半稳定希格斯丛，它没有形变），并使规范对称性保持完全未破缺. 因此，在亏格\footnote{这里以及下文关于亏格为 1 的讨论中，与本文始终采用的情形一样，我们考虑几何朗兰兹的无分歧\index{unramified@无分歧}情形. 分歧点\index{ramification-point@分歧点}将约化规范对称性，如 \cite{GW} 中所讨论的.}为 0 的情况下，膜的描述将基于一个规范群为 $G$ 的低能规范理论，而没有伴随的 $\sigma$-模型. 对于亏格为 1 的情形，在丛 $E$ 拓扑平凡时，一般的稳定希格斯丛保留 $G$ 的一个与其极大环面同构的未破缺子群. 因此，在亏格为 1 时，单非阿贝尔群 $G$ 的情形有些类似于 $G=U(1)$ 的情形：膜的一般描述将使用一个与 $\sigma$-模型相结合的阿贝尔规范理论. 对于亏格 $>1$ 的情形，一般的膜可以用 $\sigma$-模型的语言来描述，但对于支撑在 $\MH$ 奇点上的膜，这种描述就会失效，因为那里存在增强的规范对称性. （增强的规范对称性为有限群的情形已在 \cite{FW} 中研究过.）

回到 $G=U(1)$，磁通量为 $d$ 的膜的 $S$-对偶，在经典层面上可以描述为另一类新颖的膜，我们称之为电通量为 $d$ 的膜. 对于这种膜，$\phi_0$ 仍然是 $x^1$ 的线性函数，但现在存在的是恒定电场，而不是恒定磁场. 为保持超对称性，这些场必须满足适当的超对称方程（\cite{KW} 中的式（3.29），应取 $t=i$，这是研究威尔逊线的适当取值）. 在阿贝尔规范理论中，这些方程在欧几里得号差下约化为 $F+i\,d\phi=0$. 在四维流形 $M=\Bbb{R}^2\times C$ 上，以 $x^0,$ $x^1$ 为 $\Bbb{R}^2$ 上的欧几里得坐标，可以令 $\phi_0$ 为 $x^1$ 的线性函数、$F_{01}$ 为常数，并让其他一切都消失，从而解出这些方程. 在欧几里得号差下，$F_{01}$ 必须为虚数，但在洛伦兹号差下，$i$ 因子就会消失；而洛伦兹号差当然才是物理学真正的所在.

在这一分析中起重要作用的场 $\phi_0$，属于包含 $U(1)$ 规范场的超对称多重态，称为向量多重态\index{vector-multiplet@向量多重态}. 携带电荷和磁荷的膜，可以在由 $\sigma$-模型与 $U(1)$ 向量多重态耦合而成的联合低能理论中描述.



\subsection{非阿贝尔推广}
\label{actsun}
\subsubsection{计算}

现在转向非阿贝尔推广. 为考察基本思想，取 $^LG=G=U(N)$，并取磁权 $^Lw=(1,0,0,\dots,0)$，它是 $N$ 维表示 ${\cal V}$ 的最高权.

令 ${\cal B}$ 为支撑在点 $x\in \MH$ 上的零维膜，该点对应于希格斯丛 ${\cal E}=(E,\varphi)$. 插入在点 $y\times \pp\in I\times C$ 处的威尔逊算符 $W({}^Lw;y,\pp)$ 按如下方式作用：
\begin{equation}\label{howacts} W({}^Lw;y,\pp){\cal B}={\cal
B}\otimes \EE_{H,\cal V}|_{p\times x}\end{equation} 这里，$\EE_{H,\cal V}\to C\times \MH$ 是表示 ${\cal V}$ 中的万有希格斯丛（这一概念在第 \ref{eleig} 节中引入），而 $\EE_{H,\cal V}|_{p\times x}$ 是它在 $\pp\times x\in C\times \MH$ 处的限制. 因此，$\EE_{H,\cal V}|_{\pp\times x}$ 是一个固定的向量空间，而式（\ref{howacts}）表达了 ${\cal B}$ 是电本征膜这一事实.

一般地选取 $\pp$ 和 ${\cal E}$，使 $\varphi(\pp)$ 是正则半单的. 因而，在 $\pp$ 附近，本征值方程存在 $N$ 个不同的解：
\begin{equation}\label{regsol}
\varphi\uppsi_i=y_i\uppsi_i,~i=1,\dots,N.\end{equation}
相应地，$\EE_{H,\cal V}|_{\pp\times x}$ 分解为 $\varphi$ 的本征空间之和，我们将这些本征空间记为 $\EE^{(i)}_{H,\cal V}|_{\pp\times x}.$ 这些本征空间具有用谱覆盖\index{spectral-cover@谱覆盖} $\psi:D\to C$ 表述的熟知解释. 希格斯丛 ${\cal E}=(E,\varphi)$ 由一个零次线丛 ${\cal M}\to D$ 确定. 位于 $\pp\in C$ 之上的点 $\qq_i\in D$，与 $\varphi(\pp)$ 的本征向量一一对应；而 ${\cal M}$ 的定义使得 $\EE^{(i)}_{H,\cal V}|_{\pp\times x}$ 就是 ${\cal M}|_{\qq_i}$，即 ${\cal M}$ 在 $\qq_i$ 处的纤维. 因此，本征膜\index{eigenbrane@本征膜}方程（\ref{howacts}）可以更明确地写为
\begin{equation}\label{owacts} W({}^Lw;y,\pp){\cal B}={\cal
B}\otimes \left(\oplus_{i=1}^N{\cal M}|_{\qq_i}\right).\end{equation}

与 ${\cal B}$ 对偶的磁本征膜 $\hat {\cal B}$，支撑在希钦纤维化中包含点 $x$ 的纤维 $\FF$ 上. 它的陈-帕顿丛 ${\cal U}$ 由第 \ref{backdet} 节中对偶性的分析确定. 在 $\FF$ 中与零次线丛 ${\cal L}\to D$ 对应的点处，$\U$ 的纤维为 ${\cal U}|_{\cal L}=\langle{\cal L},{\cal M}\rangle$.

现在可以确定 't Hooft 算符 $T({}^Lw;y,\pp)$ 的作用. 算符 $T({}^Lw;y,\pp)$ 可以通过 ${\cal E}=(E,\varphi)$ 的 $N$ 种可能的 $\varphi$-不变赫克修正\index{hecke-modification@赫克修正}中的任意一种来作用. 它们对应于全纯映射 $\Phi_i:\FF\to \FF$，并与位于 $\pp$ 之上的点 $\qq_i\in D$ 自然地一一对应. $\Phi_i$ 按 ${\cal L}\to {\cal L}(\qq_i)$ 作用，因此把 ${\cal W}$ 映到一个新的陈-帕顿丛，该丛在 ${\cal L}$ 处的纤维为 $\langle{\cal L}(\qq_i),{\cal M}\rangle,$ 根据配对 $\langle~,~\rangle$ 的双线性，它等于 $\langle{\cal L},{\cal M}\rangle\otimes {\cal M}|_{\qq_i}$. 因此，$\Phi_i$ 把 $\hat \B$ 映到 $\hat \B\otimes {\cal M}|_{\qq_i}$；对 $\Phi_i$ 的所有选择求和后，得到
\begin{equation}\label{gladdo}T({}^Lw;y,\pp)\hat{\cal B}=\hat{\cal
B}\otimes \left(\oplus_{i=1}^N{\cal M}|_{\qq_i}\right).\end{equation} 这与式（\ref{howacts}）很好地平行对应，因此在这个例子中，我们验证了所预期的对偶性.


可以用类似的方式考虑 $^LG=U(N)$ 的其他微小权\index{minuscule-weight@微小权}. 与式（\ref{gelfix}）的讨论一样，令 $^Lw(k)=(1,1,\dots,1,0,\dots,0)$（其中 1 的个数为 $k$）为 $U(N)$ 的表示 $\wedge^k{\cal V}$ 的最高权. 首先考虑威尔逊算符 $W({}^Lw(k);y,\pp)$ 的作用. 它在零维膜 ${\cal B}$ 上按如下方式作用：
\begin{equation}\label{owactos} W({}^Lw(k);y,\pp){\cal B}={\cal
B}\otimes \EE_{H,\wedge^k\cal V}|_{\pp\times x}\end{equation}
$\EE_{H,\wedge^k\cal V}$ 是表示 $\wedge^k\cal V$ 中的万有希格斯丛，而 $\EE_{H,\wedge^k\cal V}|_{\pp\times x}$ 是它在 $\pp\times x$ 处的限制. 一旦将 ${\cal V}$ 分解为一维本征空间 ${\cal M}|_{\qq_i}$ 的直和，$\wedge^k{\cal V}$ 也就具有类似的分解 $\oplus_{\alpha} {\cal M}|_\alpha$，其中 $\alpha$ 是集合 $\{1,2,\dots,N\}$ 的一个基数为 $k$ 的子集，并且 ${\cal M}|_\alpha=\otimes_{i\in\alpha}{\cal M}|_{\qq_i}$. 因此，可以将本征值方程写为
\begin{equation}\label{turme}W({}^Lw(k);y,\pp){\cal B}={\cal
B}\otimes \left(\oplus_{\alpha}{\cal M}|_\alpha\right).\end{equation}

在磁侧，'t Hooft 算符 $T({}^Lw(k);y,\pp)$ 通过一个 $\varphi$-不变赫克变换来作用. 对每个 $\alpha$，都有一个按 ${\cal L}\to {\cal L}_\alpha={\cal L}\otimes\left(\otimes_{i\in \alpha}{\cal O}(\qq_i)\right)$ 作用的 $\varphi$-不变赫克变换 $\Phi_\alpha$. 陈-帕顿丛 $\langle{\cal L},{\cal M}\rangle$ 被 $\Phi_\alpha$ 映到 $\langle{\cal L}_\alpha,{\cal M}\rangle$；再次利用配对的双线性，后者等于 $\langle{\cal L},{\cal M}\rangle\otimes {\cal M}|_\alpha$. 因此，一旦对 $\alpha$ 求和，'t Hooft 算符就按如下方式作用：
\begin{equation}\label{urme}T({}^Lw(k);y,\pp)\hat{\cal B}=\hat{\cal
B}\otimes\left(\oplus_\alpha {\cal M}|_\alpha\right).\end{equation} 再次比较式（\ref{turme}）和（\ref{urme}），便看到了预期的对偶性.

对上述情形稍作推广，就得到 $G=U(N)$ 最一般的微小权. 引入一个整数 $r$，并令 $^Lw(k;r)=(r+1,r+1,\dots,r+1,r,\dots,r)$，其中 $k$ 个权等于 $r+1$，其余的权等于 $r$. 例如，如果 $r=-1$ 且 $k=N-1$，就得到 $^Lw=(0,0,\dots,0,-1)$，这是与 ${\cal V}$ 对偶的表示 ${\cal V}^*$ 的最高权. 一般而言，$^Lw(k;r)$ 是表示 $\wedge^k{\cal V}\otimes (\det\,{\cal V})^r$ 的最高权. 在电侧，引入 $r$ 的效果是，式（\ref{turme}）的右边必须张量乘以 $(\det\,\EE_{H,{\cal V}}|_{\pp\times x})^r$，它等于 $(\otimes_{i=1}^N{\cal M}|_{\qq_i})^r$：
\begin{equation}\label{zurme}W({}^Lw(k;r);y,\pp){\cal B}={\cal
B}\otimes \left(\oplus_{\alpha}{\cal
M}|_\alpha\right)\otimes\left(\otimes_{i=1}^N{\cal
M}|_{\qq_i}\right)^r.\end{equation} 在磁侧，$\varphi$-不变赫克修正现在按 ${\cal L}\to{\cal L}_\alpha\otimes(\otimes_{i=1}^N{\cal O}(\qq_i))^r$ 作用. 重复式（\ref{urme}）的推导，现在在磁侧得到
\begin{equation}\label{ozurme}T({}^Lw(k;r))\hat{\cal B}=\hat{\cal
B}\otimes\left(\oplus_{\alpha}{\cal
M}|_\alpha\right)\otimes\left(\otimes_{i=1}^N{\cal
M}|_{\qq_i}\right)^r.\end{equation}
这再次展示了对偶性.

\subsubsection{解释}\label{interp}

$^Lw(k;r)$ 的各个权之和为 $k+rN$；它可以为正，也可以为负. 考虑用由微小权 $^Lw(k_\sigma;r_\sigma)$、$\sigma=1,\dots,s$ 确定的一组 't Hooft 算符的乘积来作用. 这一 't Hooft 算符乘积使 $c_1(E)$ 改变 $\Delta=\sum_\sigma\left(k_\sigma+Nr_\sigma\right)$. 现在回顾一下，虽然 $G=U(N)$ 的希格斯丛 ${\cal E}=(E,\varphi)$ 可以具有任意取值的 $c_1(E)$，但作为 $\sigma$ 模型靶空间的超凯勒\index{hyperkahler@超凯勒}流形 $\MH$，恰好只参数化 $c_1(E)=0$ 的希格斯丛\transnote{$k+rN$ 是权 $^Lw(k;r)=(r+1,\dots,r+1,r,\dots,r)$ 之和，即表示 $\wedge^k{\cal V}\otimes(\det{\cal V})^r$ 的权和；'t Hooft 算符按此权平移 $c_1(E)$，故 $\Delta$ 就是移位量。取 $^L\neg G=SU(N)$（对应 $G=SU(N)/\Z_N$）时权在 $\Z_N$ 上的和为零，$\Delta$ 只能是 $N$ 的倍数，而 $PSU(N)$ 描述中的 $c_1$ 本就只定到模 $N$，于是不必额外保留 $U(1)$ 多重态——这正是本节末改用 $SU(N)$ 的理由。[来源：原文 §7.2.2；2026-11-19 核对。]}.

如果 $\Delta=0$，给定的 't Hooft 算符乘积就会使我们仍然留在这一类之内. 否则，就会遇到第 \ref{elmag} 节在 $U(1)$ 规范理论中讨论过的同一现象. 对于 $G=U(N)$，一般的希格斯丛具有 $U(1)$ 对称性，原因很简单：$U(N)$ 的中心就是 $U(1)$. 因此，低能理论是靶空间为 $\MH$ 的 $\sigma$-模型与超对称 $U(1)$ 规范理论的乘积. 与第 \ref{elmag} 节一样，纳入 $U(1)$ 理论后，就存在带有非零电荷或磁荷的膜. 在 $G=U(N)$ 的语境中，对于这些膜，$\Tr\,\phi_0$ 是 $x^1$ 的线性函数.

满足 $\Delta\not=0$ 的一组 't Hooft 算符的乘积，会平移膜的磁荷，并将一个可以完全在 $\sigma$-模型中描述的膜，映到一个其描述确实要求我们纳入 $U(1)$ 向量多重态的膜. 同样，满足 $\Delta\not=0$ 的一组威尔逊算符的乘积，会平移膜的电荷. 与以往一样，带电的膜的陈-帕顿丛，是 $\MH$ 上的扭曲线丛（或扭曲向量丛），而不是普通的线丛（或向量丛）. 这在上述分析中的体现是：威尔逊算符作用在零维膜 ${\cal B}$ 上时的“本征值”，依赖于线丛 ${\cal M}$，而不仅仅依赖于其同构类. 如果用满足 $\Delta=0$ 的威尔逊算符乘积作用，这种依赖就会相互抵消. 与通常一样，陈-帕顿波函数依赖于万有丛选择（在当前语境中，这意味着选择 ${\cal M}$）的膜，就是携带电荷的膜. 但对于中心维数为正的群 $U(N)$ 或 $U(1)$，要恰当地描述携带电荷或磁荷的膜，必须将 $U(1)$ 向量多重态纳入描述.

这些问题完全源于 $U(N)$ 的中心 $U(1)$ 并非有限阶. 因此，如果取 $^LG=SU(N)$ 而不是 $U(N)$，就可以在保留上述几何朗兰兹对偶性论证精神的同时，避免这些问题. 当然，这对应于 $G=SU(N)/\Z_N$. 对 $SU(N)$ 的微小表示\index{minuscule-representation@微小表示}，可以完全按照我们对 $U(N)$ 所用的方式来分析几何朗兰兹对偶性.


\vskip.5cm
本研究得到美国国家科学基金会项目 PHY-1606531 的部分资助.
